\subsection{Transpose and Symmetry}

\begin{definitionbox}[Transpose]
For $A=(a_{ij})\in\mathbb R^{m\times n}$, the transpose is the
$n\times m$ matrix
\[
A^{\mathsf T}=(a_{ji}).
\]
Rows of $A$ become columns of $A^{\mathsf T}$, and columns become rows.
\end{definitionbox}

\begin{examplebox}[Transposing a rectangular matrix]
\[
\begin{pmatrix}
1&2&3\\
4&5&6
\end{pmatrix}^{\mathsf T}
=
\begin{pmatrix}
1&4\\
2&5\\
3&6
\end{pmatrix}.
\]
The original matrix has size $2\times3$, while its transpose has size
$3\times2$.
\end{examplebox}

\begin{definitionbox}[Symmetric matrix]
A square matrix is symmetric when
\[
A^{\mathsf T}=A.
\]
Equivalently, its entries satisfy $a_{ij}=a_{ji}$, so the matrix is mirrored
across the main diagonal.
\end{definitionbox}

\begin{interpretationbox}[What symmetry makes visible]
In a symmetric matrix, information recorded in position $(i,j)$ agrees with the
information in position $(j,i)$. The pattern can therefore be recognised before
any calculation is performed.
\end{interpretationbox}
